\section{Arquitectura MAE-TMS}
\label{sec:arquitectura}

Esta sección describe la arquitectura completa del sistema de memoria transactiva
multi-agente.
Se presentan primero los componentes (Sección~\ref{sec:componentes}), luego los tres
diagramas de flujo que ilustran el llenado de memorias, la fase temprana y la fase
madura (Sección~\ref{sec:diagramas}), y finalmente la formalización de todas las
operaciones del sistema (Sección~\ref{sec:operaciones}).

\subsection{Componentes del sistema}
\label{sec:componentes}

El sistema MAE-TMS contiene en total \textbf{14 memorias asociativas} distribuidas entre
tres agentes especialistas y el mediador TME.
La Tabla~\ref{tab:memorias} resume todas las memorias, sus dominios y sus dimensiones.
Fuera del sistema de memorias, dos componentes de percepción actúan como interfaz con
el mundo: un encoder ResNet18 \cite{he2016resnet} que proyecta imágenes de $128\times128$
píxeles al espacio latente $\mathbf{z} \in \mathbb{R}^{64}$, y un decoder convolucional
transpuesto que invierte esa proyección; ambos son preentrenados y sus pesos permanecen
fijos durante la operación del sistema de memorias.
Las representaciones semánticas se construyen con embeddings fastText 300D
\cite{bojanowski2017fasttext} procesados por spaCy \cite{honnibal2020spacy} y
enriquecidos con relaciones de ConceptNet 5.7 \cite{speer2017conceptnet}.

\begin{table}[H]
\centering
\caption{Las 14 memorias asociativas del sistema MAE-TMS.
Para las memorias homo-asociativas se indica $(n, m)$ (dimensión, alfabeto).
Para las hetero-asociativas se indica $(n_L, m_L \leftrightarrow n_R, m_R)$ o
$(n_L, m_L \to n_R, m_R)$ cuando la relación es unidireccional.
El dominio \emph{label} corresponde a vectores fastText cuantizados;
el dominio \emph{latente} a códigos ResNet18 cuantizados.
$g \in \{0,1,2\}$ codifica la identidad del agente ganador en one-hot con bit de
\emph{desconocido}.}
\label{tab:memorias}
\small
\begin{tabular}{@{}llll@{}}
\toprule
\textbf{Componente} & \textbf{Memoria} & \textbf{Tipo} & \textbf{Dimensiones} \\
\midrule
\multirow{4}{*}{Agente $k$ (×3)} &
  $M^L_{\mathrm{dom}}$ & Homo-asociativa & $(n{=}300,\; m{=}16)$ \\
& $M^R_{\mathrm{dom}}$ & Homo-asociativa & $(n{=}64,\; m{=}32)$ \\
& $M^H_{\mathrm{dom}}$ & Hetero-asociativa & $(300, 16 \leftrightarrow 64, 32)$ \\
& $M_{\mathrm{dir}}$   & Hetero-asociativa & $(300, 16 \to 3, 2)$ \\
\midrule
\multirow{2}{*}{TME} &
  $M^L_{\mathrm{dir}}$ & Hetero-asociativa & $(300, 16 \to 3, 2)$ \\
& $M^R_{\mathrm{dir}}$ & Hetero-asociativa & $(64, 32 \to 3, 2)$ \\
\bottomrule
\end{tabular}
\end{table}

Las memorias de directorio ($M_{\mathrm{dir}}$, $M^L_{\mathrm{dir}}$,
$M^R_{\mathrm{dir}}$) son hetero-asociativas cuyo dominio derecho es la identidad del
agente ganador codificada en \emph{one-hot} con $q=2$ (un bit de ``conoce'' y uno de
``no-conoce'').
Cuando se registra un par $(\mathbf{v}_q, g)$ en un directorio, la representación
semántica (o latente) del cue se asocia al identificador del agente que ganó el
enrutamiento; todos los agentes, incluidos los perdedores, registran esta asociación
(\emph{directory updating} de Wegner).

\subsection{Diagramas de flujo}
\label{sec:diagramas}

Los tres diagramas siguientes ilustran las tres etapas operativas del sistema:
el llenado de las memorias de dominio, la recuperación mediada por el TME (fase temprana)
y la recuperación punto a punto (fase madura).

% ------------------------------------------------------------------
% Diagrama 1: Llenado
% ------------------------------------------------------------------
\begin{figure}[H]
\centering
\begin{tikzpicture}[
    node distance = 0.55cm and 1.5cm,
    box/.style = {
        rectangle, draw, line width=0.4pt,
        align=center, font=\small,
        minimum height=1.0cm, minimum width=2.6cm,
        inner sep=4pt
    },
    boxg/.style = {
        rectangle, draw=gray, dashed, line width=0.4pt,
        align=center, font=\small,
        minimum height=1.0cm, minimum width=2.6cm,
        inner sep=4pt, text=gray
    },
    arr/.style  = {-{Stealth[length=2mm]}, line width=0.4pt},
    lbl/.style  = {font=\footnotesize, midway, above, align=center},
    lblb/.style = {font=\footnotesize, midway, below, align=center},
]

% ---- Rama semántica (izquierda) ----
\node[box] (conceptnet) {ConceptNet 5.7\\relaciones \textit{IsA},\\ \textit{HasProp.}, \textit{Related}};
\node[box, right=of conceptnet] (ft) {fastText\\subword 300D};
\node[box, right=of ft] (sign) {$\mathrm{sign}(\cdot)$};
\node[box, right=of sign] (qL) {Cuantización\\binaria};

% ---- Memoria label ----
\node[box, right=1.8cm of qL] (ML) {$M^L_{\mathrm{dom}}$\\homo\\$(300,16)$};

% ---- Etiqueta secuencia ----
\node[box, below=1.5cm of qL] (seq) {Secuencia\\de labels\\(expandida por\\frecuencia)};

% ---- Rama visual (debajo de rama semántica) ----
\node[box, below=3.5cm of conceptnet] (eth80) {ETH-80\\200 imágenes/clase};
\node[box, right=of eth80] (resnet) {Encoder\\ResNet18\\$128\!\times\!128$};
\node[box, right=of resnet] (zR) {$\mathbf{z}\in\mathbb{R}^{64}$};
\node[box, right=of zR] (qR) {Cuantización\\global\\min--max};

% ---- Memoria latente y het ----
\node[box, right=1.8cm of qR] (MR) {$M^R_{\mathrm{dom}}$\\homo\\$(64,32)$};
\node[box, below=0.7cm of ML]  (MH) {$M^H_{\mathrm{dom}}$\\hetero\\$(300,16\!\leftrightarrow\!64,32)$};

% ---- Flechas rama semántica ----
\draw[arr] (conceptnet) -- node[lbl]{palabras,\\frec.} (ft);
\draw[arr] (ft) -- node[lbl]{$\mathbf{v}\in\mathbb{R}^{300}$} (sign);
\draw[arr] (sign) -- node[lbl]{$\{-1,1\}^{300}$} (qL);
\draw[arr] (qL) -- node[lbl]{$\mathbf{v}_q\in\{0,15\}^{300}$} (ML);
\draw[arr] (qL) -- (seq);

% ---- Flechas rama visual ----
\draw[arr] (eth80) -- (resnet);
\draw[arr] (resnet) -- node[lbl]{$\mathbf{z}$} (zR);
\draw[arr] (zR) -- (qR);
\draw[arr] (qR) -- node[lbl]{$\mathbf{z}_q\in[0,31]^{64}$} (MR);

% ---- Registro emparejado image-major en M^H ----
\draw[arr] (qL.south) -- ++(0,-0.4) -| (MH.north)
    node[near end, left, font=\footnotesize]{$\lambda(\mathbf{v}_q, \mathbf{z}_q)$};
\draw[arr] (qR.north) |- (MH.west);
\draw[arr] (seq.east) -- ++(0.4,0) |- (MH.south west);

% ---- Nota ----
\node[font=\footnotesize\itshape, text width=4.5cm, align=center,
      below=0.3cm of MH]
    {La distribución de clase se acumula DENTRO\\
     de la memoria; el prototipo emerge, no se impone.};

\end{tikzpicture}
\caption{Llenado de las memorias de dominio (etapas 3--5).
La rama semántica extrae y cuantiza embeddings fastText desde ConceptNet;
la rama visual cuantiza los códigos ResNet18.
El registro en $M^H_{\mathrm{dom}}$ sigue el esquema \emph{image-major}:
la imagen $i$-ésima se registra una sola vez, emparejada con el label
$\mathbf{v}_q^{(i mod L)}$ de la secuencia expandida por frecuencias
(de longitud $L$), de modo que cada clase aporta exactamente $N=200$ registros.}
\label{fig:arq-llenado}
\end{figure}

% ------------------------------------------------------------------
% Diagrama 2: Fase temprana
% ------------------------------------------------------------------
\begin{figure}[H]
\centering
\begin{tikzpicture}[
    node distance = 0.6cm and 1.4cm,
    box/.style = {
        rectangle, draw, line width=0.4pt,
        align=center, font=\small,
        minimum height=1.0cm, minimum width=2.5cm,
        inner sep=4pt
    },
    boxA/.style = {
        rectangle, draw, line width=0.4pt,
        align=center, font=\small,
        minimum height=2.2cm, minimum width=2.8cm,
        inner sep=4pt
    },
    arr/.style  = {-{Stealth[length=2mm]}, line width=0.4pt},
    lbl/.style  = {font=\footnotesize, midway, above, align=center},
    lblb/.style = {font=\footnotesize, midway, below, align=center},
]

% ---- Entrada ----
\node[box] (query) {Consulta};
\node[box, right=of query] (spacy) {spaCy\\lemas NOUN/ADJ};
\node[box, right=of spacy] (emb) {fastText\\$+\,\mathrm{sign}$};
\node[box, right=of emb] (vq) {$\mathbf{v}_q$};

% ---- TME ----
\node[box, right=1.8cm of vq, minimum height=1.4cm, minimum width=2.8cm]
    (tme) {TME\\(mediador)};

% ---- Agentes (col derecha) ----
\node[boxA, right=2.5cm of tme, yshift=2.0cm]  (A0) {Agente 0\\(\textit{apple})\\$M^L \to w \to M^H$};
\node[boxA, right=2.5cm of tme]                 (A1) {Agente 1\\(\textit{horse})\\$M^L \to w \to M^H$};
\node[boxA, right=2.5cm of tme, yshift=-2.0cm]  (A2) {Agente 2\\(\textit{car})\\$M^L \to w \to M^H$};

% ---- Scores ----
\node[box, right=1.2cm of A1] (argmax) {$\arg\max_k \sum s_k$\\ganador $g$};

% ---- Salidas ----
\node[box, right=of argmax, yshift=1.2cm] (recall) {Recall $\beta$\\$\hat{\mathbf{z}}_q$};
\node[box, right=of recall] (decoder) {Decoder};

% Corchete de salida
\node[box, right=of decoder] (salida) {Imagen\\recuperada};

% ---- Actualización de directorio ----
\node[box, below=1.4cm of tme, minimum width=2.8cm]
    (dir_update) {Actualización de\\directorio\\ $(\mathbf{v}_q \to g)$};

% ---- Flechas principales ----
\draw[arr] (query) -- (spacy);
\draw[arr] (spacy) -- (emb);
\draw[arr] (emb) -- node[lbl]{$\mathbf{v}\in\mathbb{R}^{300}$} (vq);
\draw[arr] (vq) -- (tme);

% Broadcast del TME a los agentes (3 pares de flechas paralelas)
\foreach \ag/\lab in {A0/$s_0$, A1/$s_1$, A2/$s_2$}{
    \draw[arr] (tme.east) -- (\ag.west);
    \draw[arr] (\ag.east) -- node[lbl]{\lab} (argmax.west|-\ag.east);
}

\draw[arr] (argmax) -- (recall);
\draw[arr] (recall) -- (decoder);
\draw[arr] (decoder) -- (salida);

% Actualización de directorio hacia los 3 agentes y el TME
\draw[arr] (tme.south) -- (dir_update.north);
\draw[arr, dashed] (dir_update.west) -| (A0.south);
\draw[arr, dashed] (dir_update)      -- (A1.south);
\draw[arr, dashed] (dir_update.east) -| (A2.south);
\node[font=\footnotesize, below=0.15cm of dir_update, text width=5cm, align=center]
    {Todos los agentes y el TME\\ registran $\lambda(\mathbf{v}_q, g)$};

% Nota Wegner
\node[font=\footnotesize\itshape, below=0.25cm of A2, text width=4cm, align=center]
    {También los perdedores anotan\\(\textit{directory updating}, Wegner)};

\end{tikzpicture}
\caption{Fase temprana: aprendizaje mediado por el TME.
La consulta se codifica con fastText+sign, el TME la difunde a los tres agentes,
cada agente devuelve un score $s_k$ mediante el reconocimiento de su memoria homo
y la proyección de su memoria hetero, y el ganador $g = \arg\max_k \sum s_k$ ejecuta
el recall $\beta$ hacia el espacio latente.
Los tres agentes y el propio TME registran el par $(\mathbf{v}_q, g)$ en sus
respectivos directorios, incluyendo los agentes perdedores.}
\label{fig:arq-temprana}
\end{figure}

% ------------------------------------------------------------------
% Diagrama 3: Fase madura
% ------------------------------------------------------------------
\begin{figure}[H]
\centering
\begin{tikzpicture}[
    node distance = 0.6cm and 1.3cm,
    box/.style = {
        rectangle, draw, line width=0.4pt,
        align=center, font=\small,
        minimum height=1.0cm, minimum width=2.6cm,
        inner sep=4pt
    },
    boxg/.style = {
        rectangle, draw=gray, dashed, line width=0.4pt,
        align=center, font=\small, text=gray,
        minimum height=1.2cm, minimum width=2.6cm,
        inner sep=4pt
    },
    arr/.style  = {-{Stealth[length=2mm]}, line width=0.4pt},
    arrg/.style = {-{Stealth[length=2mm]}, line width=0.4pt, draw=gray, dashed},
    lbl/.style  = {font=\footnotesize, midway, above, align=center},
]

% ---- Entrada ----
\node[box] (query) {Consulta};
\node[box, right=of query] (entrada) {Agente de\\entrada\\(aleatorio)};

% ---- Directorio del agente de entrada ----
\node[box, right=1.6cm of entrada] (mdir) {$M_{\mathrm{dir}}$\\lectura B1};

% ---- Rama conocida: agente destino ----
\node[box, right=2.0cm of mdir, yshift=1.4cm] (destino) {Agente\\destino};
\node[box, right=of destino] (recall) {Recall $\beta$\\$\hat{\mathbf{z}}_q$};
\node[box, right=of recall] (decoder) {Decoder};
\node[box, right=of decoder] (img) {Imagen};

% ---- Rama desconocida: rechazo ----
\node[box, right=2.0cm of mdir, yshift=-1.4cm] (rechazo) {RECHAZO\\(ningún experto)};

% ---- TME (fuera del circuito) ----
\node[boxg, below=1.5cm of entrada] (tme) {TME\\(fuera del circuito)};

% ---- Flechas ----
\draw[arr] (query) -- (entrada);
\draw[arr] (entrada) -- node[lbl]{$\mathbf{v}_q$} (mdir);

% Bifurcación desde el directorio
\draw[arr] (mdir.east) -- ++(0.6,0) |- (destino.west)
    node[near start, above, font=\footnotesize]{$g$ conocido};
\draw[arr] (mdir.east) -- ++(0.6,0) |- (rechazo.west)
    node[near start, below, font=\footnotesize]{$c{+}\text{desconocido}$};

% Agente destino al recall
\draw[arr] (destino) -- (recall);
\draw[arr] (recall) -- (decoder);
\draw[arr] (decoder) -- (img);

% TME apagado
\draw[arrg] (entrada.south) -- (tme.north);
\node[font=\footnotesize\itshape, text=gray, below=0.2cm of tme, text width=4cm, align=center]
    {El mediador ya no participa\\en la recuperación};

% Nota rechazo
\node[font=\footnotesize\itshape, right=0.3cm of rechazo, text width=3.5cm, align=left]
    {El grupo no inventa\\expertos inexistentes};

% Corchete salida final
\node[right=0.15cm of img, font=\small] (cbrace) {$\left.\rule{0pt}{2.5ex}\right\}$};
\node[right=0.1cm of cbrace, font=\footnotesize, align=left]
    {imagen,\\labels,\\ubicación};

\end{tikzpicture}
\caption{Fase madura: recuperación punto a punto.
Un agente de entrada (seleccionado aleatoriamente) consulta su directorio $M_{\mathrm{dir}}$
con la lectura B1 y reenvía directamente al agente destino sin pasar por el TME
(dibujado en gris punteado).
Si el cue cae en la región ``desconocido'' del directorio, el sistema emite un rechazo
explícito: el grupo no inventa expertos.}
\label{fig:arq-madura}
\end{figure}

\subsection{Operaciones del sistema}
\label{sec:operaciones}

\subsubsection*{Score de enrutamiento (gated)}

Para un agente $k$ dado un cue $\mathbf{v}_q$ (representación semántica cuantizada),
el proceso de enrutamiento es el siguiente.

\begin{enumerate}
    \item Se obtienen los pesos de reconocimiento de la memoria homo-asociativa de labels:
    \begin{equation}
        \mathbf{w} = \eta(\mathbf{v}_q,\; M^L_{\mathrm{dom}})
        \label{eq:pesos-routing}
    \end{equation}
    y se normalizan $w_i \leftarrow w_i / \max_j w_j$ para equiparar escalas.

    \item Se computa la proyección AND-ponderada sobre la memoria de contenido:
    \begin{equation}
        \Pi_k = \Pi(\mathbf{v}_q,\; \mathbf{w},\; M^H_{\mathrm{dom}})
        \label{eq:proyeccion-routing}
    \end{equation}
    según la ecuación~\eqref{eq:proyeccion-AND}.

    \item El score del agente $k$ incorpora un \textbf{gate de containment}:
    \begin{equation}
        s_k = \begin{cases}
            0 & \text{si } \exists\, j : \sum_l \Pi_k[j, l] = 0 \\[4pt]
            \dfrac{\displaystyle\sum_{j,l} \Pi_k[j, l]}%
                  {|\{(j,l) : \Pi_k[j,l] \neq 0\}|} & \text{en otro caso}
        \end{cases}
        \label{eq:score-gated}
    \end{equation}
    El gate anula el score si alguna fila del dominio derecho carece completamente
    de soporte: el agente no opina sobre lo que no contiene.
    En el caso contrario, $s_k$ es la \emph{activación media} de las celdas no nulas.
\end{enumerate}

El agente ganador es $g = \arg\max_k \sum_{\text{tokens}} s_k$ (la suma recorre todos
los tokens semánticos del cue).
Si $s_k = 0$ para todos los agentes, el sistema emite un \textbf{rechazo}.

\paragraph{Justificación.}
Con el llenado \emph{image-major}, cada agente acumula exactamente 200 registros en
$M^H_{\mathrm{dom}}$ (una imagen por cada label de su clase, expandida por frecuencias).
Las masas de registro quedan igualadas por construcción, de modo que la activación media
$s_k$ es directamente comparable entre agentes sin necesidad de calibración adicional.

\subsubsection*{Lectura B1 del directorio}

Sea $\Pi^{\mathrm{dir}}_k(\mathbf{v}_q)$ la proyección del cue $\mathbf{v}_q$
en el directorio $M_{\mathrm{dir}}$ del agente, restringida al bit $b=1$
(``el agente $k$ conoce'').
El score B1 se define como:
\begin{equation}
    \hat{s}_k = \frac{\Pi^{\mathrm{dir}}_k(\mathbf{v}_q,\; b{=}1)}{c_k + 1}
    \label{eq:b1}
\end{equation}
donde $c_k$ es el número total de registros del agente $k$ en el directorio.

\paragraph{Justificación.}
La especialización genuina produce masas desiguales entre agentes: por ejemplo,
$c_0 = 68$, $c_1 = 54$, $c_2 = 42$ en una sesión típica.
El score crudo de una proyección hetero-asociativa crece con la masa de registros,
por lo que un agente con más registros en el directorio obtendría ventaja sin mérito.
La división por $c_k + 1$ convierte \emph{evidencia total} en \emph{evidencia por
registro}, haciendo que el score refleje calidad de especialización y no cantidad.
Esta es la \textbf{única calibración externa} del sistema de directorios.

\subsubsection*{Entropía del directorio}

Sea $c_k$ el número de registros del agente $k$ en el directorio y
$C = \sum_k c_k$ el total.
La entropía del directorio mide la salud estructural del grupo:
\begin{equation}
    H_{\mathrm{dir}} = -\sum_{k=0}^{K-1} p_k \log_2 p_k, \qquad p_k = \frac{c_k}{C}
    \label{eq:entropia-dir}
\end{equation}
Con $K=3$ agentes, el máximo teórico es $\log_2 3 \approx 1{.}585$ bits,
alcanzado cuando la especialización es perfectamente balanceada.
Una entropía menor señala o bien que un agente domina el directorio (colapso) o bien
que el directorio está poco poblado (formación incompleta).

\subsubsection*{Evocación inversa (imagen $\to$ labels)}

Dado un código latente cuantizado $\mathbf{z}_q$ de una imagen, la recuperación
de labels procede como sigue.

\begin{enumerate}
    \item Se obtienen los pesos del dominio latente:
    $\mathbf{w}_R = \eta(\mathbf{z}_q,\; M^R_{\mathrm{dom}})$.

    \item Se aplica el recall $\beta$ de la memoria hetero-asociativa \emph{desde la
    derecha}: el cue es $\mathbf{z}_q$ con pesos $\mathbf{w}_R$, y el resultado es un
    patrón muestreado en el dominio semántico $\hat{\mathbf{v}}_q$.

    \item Se buscan los 3 vecinos más cercanos de $\hat{\mathbf{v}}_q$ en el diccionario
    fastText mediante similitud coseno, y se devuelven como labels asociados a la imagen.
\end{enumerate}

\paragraph{Justificación del diccionario.}
El recall $\beta$ devuelve una \emph{muestra} de la distribución de labels acumulada
por el agente, no una palabra exacta: el vector muestreado puede no coincidir con ningún
embedding almacenado.
El diccionario actúa como un decodificador que interpreta esa muestra en términos del
vocabulario conocido.
La necesidad del diccionario no es una debilidad del sistema sino una consecuencia
directa de la naturaleza entrópica del recall.

\subsubsection*{Fidelidad de la fase madura}

La \textbf{fidelidad} se define como la fracción de consultas en la fase madura en las
que el agente de entrada enruta hacia el mismo agente ganador que la fase temprana:
\begin{equation}
    F = \frac{|\{q : g^{\mathrm{maduro}}(q) = g^{\mathrm{temprano}}(q)\}|}{N_{\mathrm{consultas}}}
    \label{eq:fidelidad}
\end{equation}
La fidelidad mide \emph{coincidencia con lo aprendido}, no acierto absoluto:
un sistema con fidelidad alta ha trasladado fielmente el conocimiento de la fase temprana
a los directorios individuales.
